\chapter{Conway Span-Monads and Conway Mealy Morphisms}
\label{ch:conway-span-monads-mealy}

\section{Purpose of the chapter}
\label{sec:conway-purpose}

In the previous chapter we studied the bicategory
\[
\mathbf{Span}(\mathcal E)
\]
and used it to identify internal categories with monads in spans.  We also
saw that the morphisms naturally induced by span-monad morphisms are Mealy
morphisms.  In the notation fixed there, a Mealy morphism
\[
(P,\Phi):\mathbb A\rightsquigarrow\mathbb B
\]
is carried by a span
\[
A_0\xleftarrow{p}P\xrightarrow{q}B_0
\]
and has structure cell
\[
\Phi:P\circ A\Rightarrow B\circ P.
\]
Equivalently, it consists of transition and output maps
\[
\delta:A_1\times_{t_A,A_0,p}P\to P,
\qquad
\omega:A_1\times_{t_A,A_0,p}P\to B_1.
\]
Thus an input arrow
\[
a:u\to v
\]
is admissible at a state \(x\in P\) when
\[
p(x)=v,
\]
and the transition has the form
\[
x
\xrightarrow{\;a/\omega(a,x)\;}
\delta(a,x),
\]
where
\[
p(\delta(a,x))=u,
\qquad
\omega(a,x):q(\delta(a,x))\to q(x).
\]
The machine therefore processes input arrows target-to-source.  This
convention is used throughout the present chapter.

The goal of this chapter is to add Conway fixed-point structure to that
picture.  The first part explains the analogy with the ordinary passage from
algebraic theories to monads.  The second part applies the span-Mealy calculus
to internal Conway categories.

The guiding object-level construction is
\[
\boxed{
\text{Conway span-monad}
=
\text{span-monad together with chosen finite-product and Conway-dagger}
\text{ structure on its associated internal category.}
}
\]
The corresponding stateful morphisms are Conway Mealy morphisms.  Ordinary
internal Conway functors occur among the representable, stateless cases:
precisely, they are the representable cases whose induced internal functors
strictly preserve the chosen finite-product structure and the Conway dagger.

\section{From algebraic theories to monads}
\label{sec:conway-algebraic-theories-monads}

We begin with the ordinary algebraic pattern.

A Lawvere theory is a small category \(\mathbb T\) with chosen finite products
and a distinguished object \(G\) such that every object is a finite power
\(G^n\).  A morphism
\[
n\to m
\]
is an \(m\)-tuple of \(n\)-ary operations, and composition is substitution of
terms into terms.

The standard construction associates to \(\mathbb T\) a finitary monad
\[
T_{\mathbb T}:\mathbf{Set}\to\mathbf{Set}.
\]
The set \(T_{\mathbb T}(X)\) is the set of \(\mathbb T\)-terms in variables
from \(X\), modulo the equations of the theory.  Equivalently,
\[
T_{\mathbb T}(X)
\cong
\int^{n\in\mathbf{FinSet}}
\mathbb T(n,1)\times X^n.
\]
The unit inserts variables as terms, and the multiplication is induced by
substitution.

Conversely, if \(T\) is a finitary monad on \(\mathbf{Set}\), then its
associated Lawvere theory has finite cardinals as objects and hom-sets
\[
\mathbb T_T(n,m)
=
\mathbf{Set}(m,Tn).
\]
Equivalently,
\[
\mathbb T_T
=
\bigl(\mathbf{Kl}(T)_{\mathrm{fin}}\bigr)^{\op},
\]
where \(\mathbf{Kl}(T)_{\mathrm{fin}}\) is the full subcategory of the
Kleisli category on finite cardinals.  The finite products in
\(\mathbb T_T\) are induced by finite coproducts in the Kleisli category.

Thus one has the classical correspondence
\[
\boxed{
\{\text{Lawvere theories}\}
\simeq
\{\text{finitary monads on }\mathbf{Set}\}.
}
\]

This is the first instance of the pattern used in this chapter: algebraic
structure can be transported across a representation theorem.  For Lawvere
theories, the representation theorem is the correspondence with finitary
monads.  For the span calculus, the relevant representation theorem is the
correspondence between span-monads and internal categories.

\section{Conway theories and finitary iteration}
\label{sec:conway-theories-finitary-iteration}

A Conway theory is a Lawvere theory equipped with a coherent operation for
solving finite systems of recursive equations.  In Lawvere-theoretic notation,
a system of \(m\) recursive equations with \(n\) parameters is a morphism
\[
f:n+m\to m.
\]
The Conway operation assigns to it a chosen solution
\[
f^\dagger:n\to m.
\]

\begin{definition}[Conway theory]
\label{def:conway-theory}
A \emph{Conway theory} is a Lawvere theory \(\mathbb T\) equipped with
operations
\[
(-)^\dagger:
\mathbb T(n+m,m)\to \mathbb T(n,m)
\]
for all finite cardinals \(n,m\), satisfying the Conway fixed-point,
parameter, sliding, and double-dagger identities.
\end{definition}

Under the Lawvere-theory/finitary-monad correspondence, the same structure
appears as an iteration operation on the associated finitary monad.

Indeed, using
\[
\mathbb T(n,m)
\cong
\mathbf{Set}(m,T_{\mathbb T}n),
\]
a morphism
\[
f:n+m\to m
\]
corresponds to a function
\[
\bar f:m\to T_{\mathbb T}(n+m).
\]
Thus a finite system of recursive equations is represented monadically as
\[
e:m\to T_{\mathbb T}(n+m),
\]
and its solution is a map
\[
e^\dagger:m\to T_{\mathbb T}n.
\]

The monadic fixed-point equation is
\[
e^\dagger
=
\mu_n\circ T[\eta_n,e^\dagger]\circ e.
\]
Equivalently, the following diagram commutes:
\[
\begin{tikzcd}[column sep=large,row sep=large]
m
  \arrow[r,"e"]
  \arrow[dr,"{e^\dagger}"']
&
T(n+m)
  \arrow[r,"{T[\eta_n,e^\dagger]}"]
&
T(Tn)
  \arrow[d,"\mu_n"]
\\
&
Tn
  \arrow[r,equal]
&
Tn .
\end{tikzcd}
\]

\begin{definition}[Finitary Conway monad]
\label{def:finitary-conway-monad}
A \emph{finitary Conway monad} on \(\mathbf{Set}\) is a finitary monad
\[
T:\mathbf{Set}\to\mathbf{Set}
\]
equipped with iteration operations
\[
(-)^\dagger:
\mathbf{Set}(m,T(n+m))
\to
\mathbf{Set}(m,Tn)
\]
for finite cardinals \(m,n\), satisfying the monadic translations of the
Conway fixed-point, parameter, sliding, and double-dagger identities.
\end{definition}

With Conway-preserving morphisms on both sides, the ordinary
Lawvere-theory/finitary-monad correspondence restricts to
\[
\boxed{
\{\text{Conway theories}\}
\simeq
\{\text{finitary Conway monads on }\mathbf{Set}\}.
}
\]

This is the ordinary algebraic model for what follows.  The rest of the
chapter replaces ordinary algebraic theories by finite-limit theories, and
replaces monads on \(\mathbf{Set}\) by monads in the bicategory of spans.

\section{Finite-limit theories and internal structure}
\label{sec:finite-limit-theories-internal-structure}

The theory of categories is not an ordinary algebraic theory.  Composition is
not a total binary operation
\[
C_1\times C_1\to C_1.
\]
It is defined only on the object of composable pairs
\[
C_1\times_{C_0}C_1.
\]
Thus the appropriate syntax is finite-limit syntax, equivalently the syntax
of essentially algebraic theories or finite-limit sketches.

A finite-limit theory has sorts, finite-limit contexts, operations between
those contexts, and equations between parallel operations.  Its models in a
finitely complete category \(\mathcal E\) interpret the specified finite-limit
contexts as actual finite limits in \(\mathcal E\).

The finite-limit theory of categories, denoted
\[
\CatTh,
\]
has sorts
\[
C_0,
\qquad
C_1,
\]
operations
\[
s,t:C_1\to C_0,
\qquad
e:C_0\to C_1,
\]
and composition
\[
m:C_1\times_{t,C_0,s}C_1\to C_1,
\]
subject to the usual source-target, identity, and associativity equations.
For every category \(\mathcal E\) with finite limits,
\[
\operatorname{Mod}_{\mathcal E}(\CatTh)
\simeq
\operatorname{Cat}(\mathcal E),
\]
the category of internal categories and internal functors in \(\mathcal E\).

This finite-limit theory is the syntactic counterpart of the span-monad
description from the previous chapter:
\[
\boxed{
\text{models of }\CatTh
\quad\leftrightarrow\quad
\text{monads in }\mathbf{Span}(\mathcal E).
}
\]

\section{The finite-limit theory of cartesian categories}
\label{sec:cartesian-categories-finite-limit-theory}

The Conway dagger is defined on categories with finite products.  We therefore
extend \(\CatTh\) to a finite-limit theory of categories with chosen finite
products.

Throughout this section, \(\mathrm{pr}_1,\mathrm{pr}_2\) denote the ordinary
projections from a product context such as \(C_0\times C_0\).  The symbols
\(\pi_1,\pi_2\) are reserved for the projection arrows of the internal
cartesian category.

\begin{definition}[The theory \(\CartCatTh\)]
\label{def:cartcatth}
The finite-limit theory
\[
\CartCatTh
\]
extends \(\CatTh\) by adjoining the following operations and equations.

First, it has a chosen terminal object, represented by a morphism
\[
\top:1\to C_0.
\]
It also has a terminal-arrow operation
\[
\tau:C_0\to C_1.
\]
Externally, \(\tau(A)\) is written
\[
!_A:A\to \top.
\]
The typing equations are
\[
s\tau=1_{C_0},
\qquad
t\tau=\top\circ {!}_{C_0},
\]
where
\[
{!}_{C_0}:C_0\to 1
\]
is the unique map to the terminal object of the ambient finite-limit theory.

Second, it has a chosen binary product operation on objects
\[
\times:C_0\times C_0\to C_0.
\]
It has projection-arrow operations
\[
\pi_1,\pi_2:C_0\times C_0\to C_1
\]
with typing equations
\[
s\pi_1=\times,
\qquad
t\pi_1=\mathrm{pr}_1,
\]
and
\[
s\pi_2=\times,
\qquad
t\pi_2=\mathrm{pr}_2.
\]
Thus, externally,
\[
\pi_1(A,B):A\times B\to A,
\qquad
\pi_2(A,B):A\times B\to B.
\]

Third, it has a pairing operation.  Let
\[
P_{\mathrm{pair}}
=
C_1\times_{s,C_0,s}C_1
\]
be the object of pairs of arrows with common source.  A generalized element is
a pair
\[
(f,g)
\]
with
\[
s(f)=s(g).
\]
The pairing operation is a morphism
\[
\langle-,-\rangle:P_{\mathrm{pair}}\to C_1
\]
with typing equations
\[
s\langle f,g\rangle=s(f),
\qquad
t\langle f,g\rangle=t(f)\times t(g).
\]
It satisfies the projection equations
\[
\pi_1(t(f),t(g))\circ\langle f,g\rangle=f,
\qquad
\pi_2(t(f),t(g))\circ\langle f,g\rangle=g.
\]

Finally, the usual uniqueness equations for terminal arrows and pairings are
imposed.  Explicitly, for every arrow
\[
h:Z\to \top
\]
one has
\[
h=!_Z,
\]
and for every arrow
\[
h:Z\to A\times B
\]
one has
\[
\left\langle
\pi_1(A,B)\circ h,
\pi_2(A,B)\circ h
\right\rangle
=
h.
\]
All these equations are equations between morphisms in the corresponding
finite-limit contexts.
\end{definition}

\begin{proposition}
\label{prop:models-cartcatth}
For every finitely complete category \(\mathcal E\), models of
\(\CartCatTh\) in \(\mathcal E\) are internal categories in \(\mathcal E\)
equipped with chosen internal finite products.  Homomorphisms of
\(\CartCatTh\)-models are internal functors preserving the chosen finite
product structure strictly.
\end{proposition}

\begin{proof}
A model of \(\CartCatTh\) interprets the underlying \(\CatTh\)-structure as an
internal category.  The additional operations interpret a chosen terminal
object, chosen terminal arrows, chosen binary products, chosen projections,
and chosen pairings.  The equations assert the corresponding finite-product
universal properties internally.

A homomorphism of models preserves all operations of the finite-limit theory.
It is therefore precisely an internal functor preserving the chosen terminal
object, terminal arrows, binary products, projections, and pairings.
\end{proof}

\section{Conway categories}
\label{sec:conway-categories}

We now recall the finite-product form of the Conway identities.  We use the
bracketing convention
\[
A\times X\times X := (A\times X)\times X.
\]

\begin{definition}[Conway category]
\label{def:conway-category}
A \emph{Conway category} is a category \(\mathcal C\) with chosen finite
products together with operations
\[
(-)^\dagger_{A,X}:
\mathcal C(A\times X,X)\to \mathcal C(A,X)
\]
for all objects \(A,X\), satisfying the following identities.

\begin{enumerate}
\item \emph{Fixed point.}  For every
\[
f:A\times X\to X,
\]
one has
\[
f^\dagger
=
f\circ \langle 1_A,f^\dagger\rangle.
\]

\item \emph{Parameter identity.}  For every
\[
h:B\to A
\]
and every
\[
f:A\times X\to X,
\]
one has
\[
\bigl(f\circ(h\times 1_X)\bigr)^\dagger
=
f^\dagger\circ h.
\]

\item \emph{Sliding/dinaturality.}  For every
\[
f:A\times X\to Y,
\qquad
g:A\times Y\to X,
\]
one has
\[
\bigl(g\circ\langle \pi_A,f\rangle\bigr)^\dagger
=
g\circ
\left\langle
1_A,
\bigl(f\circ\langle \pi_A,g\rangle\bigr)^\dagger
\right\rangle.
\]

\item \emph{Double dagger.}  For every
\[
f:(A\times X)\times X\to X,
\]
one has
\[
(f^\dagger)^\dagger
=
\bigl(f\circ\Delta_{A,X}\bigr)^\dagger,
\]
where \(f\) is first regarded as a morphism
\[
(A\times X)\times X\to X
\]
and
\[
\Delta_{A,X}:A\times X\to (A\times X)\times X
\]
is the map
\[
\Delta_{A,X}
=
\left\langle
\langle \pi_A,\pi_X\rangle,
\pi_X
\right\rangle.
\]
Equivalently, in informal tuple notation,
\[
\Delta_{A,X}(a,x)=((a,x),x).
\]
\end{enumerate}
\end{definition}

A Conway theory is the special case of a Conway category whose underlying
finite-product category is generated by finite powers of a distinguished
object.

\section{The finite-limit theory of Conway categories}
\label{sec:finite-limit-theory-conway-categories}

The definition of Conway category is finite-limit algebraic.  The objects
\(A\), \(X\), the arrow \(f:A\times X\to X\), the product \(A\times X\), and
the equation defining \(f^\dagger\) are all expressible using finite-limit
contexts and operations.

Let
\[
\ConvCatTh
\]
denote the finite-limit theory obtained from \(\CartCatTh\) by adjoining a
dagger operation.

The context of inputs to the dagger is the context of triples
\[
(A,X,f)
\]
where
\[
f:A\times X\to X.
\]
Internally this context is represented by the pullback
\[
\begin{tikzcd}[column sep=large,row sep=large]
D
  \arrow[r,"\bar f"]
  \arrow[d,"q_D"']
  \arrow[dr,phantom,"\lrcorner",very near start]
&
C_1
  \arrow[d,"{(s,t)}"]
\\
C_0\times C_0
  \arrow[r,"{\langle \times,\mathrm{pr}_2\rangle}"']
&
C_0\times C_0 .
\end{tikzcd}
\]
Thus a generalized element of \(D\) is a triple
\[
(A,X,f)
\]
with
\[
s(f)=A\times X,
\qquad
t(f)=X.
\]

Write
\[
A_D:=\mathrm{pr}_1q_D:D\to C_0,
\qquad
X_D:=\mathrm{pr}_2q_D:D\to C_0.
\]

\begin{definition}[The theory \(\ConvCatTh\)]
\label{def:convcatth}
The finite-limit theory
\[
\ConvCatTh
\]
is obtained from \(\CartCatTh\) by adjoining an operation
\[
(-)^\dagger:D\to C_1.
\]
For a generalized element
\[
(A,X,f)\in D,
\]
we write the resulting arrow as
\[
f^\dagger:A\to X.
\]
The dagger operation is required first to satisfy the typing equations
\[
s(f^\dagger)=A,
\qquad
t(f^\dagger)=X,
\]
equivalently,
\[
s\circ(-)^\dagger=A_D,
\qquad
t\circ(-)^\dagger=X_D.
\]
It is then required to satisfy the fixed-point, parameter,
sliding/dinaturality, and double-dagger identities of
Definition~\ref{def:conway-category}, interpreted as equations between
morphisms in the corresponding finite-limit contexts.
\end{definition}

\begin{definition}[Internal Conway category]
\label{def:internal-conway-category}
Let \(\mathcal E\) be finitely complete.  An \emph{internal Conway category}
in \(\mathcal E\) is a model of
\[
\ConvCatTh
\]
in \(\mathcal E\).
\end{definition}

\begin{definition}[Internal Conway functor]
\label{def:internal-conway-functor}
An \emph{internal Conway functor} is a homomorphism of
\(\ConvCatTh\)-models.  Equivalently, it is an internal functor preserving the
chosen terminal object, terminal arrows, binary products, projections,
pairings, and Conway dagger.
\end{definition}

\begin{proposition}
\label{prop:models-convcatth}
For every finitely complete category \(\mathcal E\),
\[
\operatorname{Mod}_{\mathcal E}(\ConvCatTh)
\]
is the category of internal Conway categories in \(\mathcal E\) and internal
Conway functors.
\end{proposition}

\begin{proof}
This is the interpretation theorem for the finite-limit theory
\(\ConvCatTh\).  A model interprets the underlying \(\CartCatTh\)-structure as
an internal category with chosen finite products.  The additional operation
\[
(-)^\dagger:D\to C_1
\]
interprets a Conway dagger on that internal category, and the equations are
exactly the internal Conway identities.

A homomorphism of models preserves all operations in the finite-limit theory,
hence is precisely an internal functor preserving both the finite-product
structure and the dagger.
\end{proof}

\section{Conway span-monads}
\label{sec:conway-span-monads}

We now transport the preceding finite-limit structure across the span-monad
description of internal categories.

Throughout this section, let \(\mathcal E\) be finitely complete.

\begin{definition}[Conway span-monad]
\label{def:conway-span-monad}
A \emph{Conway span-monad} in
\[
\mathbf{Span}(\mathcal E)
\]
is a monad in \(\mathbf{Span}(\mathcal E)\) whose associated internal category
is equipped with chosen internal finite products and a Conway dagger.
\end{definition}

Explicitly, a Conway span-monad consists of a span-monad
\[
C_0\xleftarrow{s}C_1\xrightarrow{t}C_0
\]
with unit
\[
e:C_0\to C_1
\]
and multiplication
\[
m:C_1\times_{t,C_0,s}C_1\to C_1,
\]
together with chosen internal finite products and a Conway dagger on the
internal category
\[
(C_0,C_1,s,t,e,m).
\]

\begin{theorem}[Conway span-monads are internal Conway categories]
\label{thm:conway-span-monads-internal-conway-categories}
Let \(\mathcal E\) be finitely complete.  To give a Conway span-monad in
\[
\mathbf{Span}(\mathcal E)
\]
is equivalently to give an internal Conway category in \(\mathcal E\).
\end{theorem}

\begin{proof}
By the previous chapter, a monad in \(\mathbf{Span}(\mathcal E)\) is an
internal category in \(\mathcal E\).  By
Proposition~\ref{prop:models-convcatth}, a model of \(\ConvCatTh\) is an
internal category equipped with chosen internal finite products and a Conway
dagger.  Therefore a span-monad equipped with this additional structure is
exactly an internal Conway category.
\end{proof}

This is an equivalence of structures on objects.  The morphisms relevant to
the span presentation are introduced below as Mealy morphisms.

\section{Product-state Mealy morphisms}
\label{sec:product-state-mealy-morphisms}

We next refine the Mealy morphisms of the previous chapter.  Let
\[
\mathbb A
\qquad\text{and}\qquad
\mathbb B
\]
be internal categories with chosen finite products, and let
\[
(P,\delta,\omega):\mathbb A\rightsquigarrow\mathbb B
\]
be a Mealy morphism, carried by
\[
A_0\xleftarrow{p}P\xrightarrow{q}B_0.
\]

A state \(x\in P\) has an input interface \(p(x)\in A_0\) and an output
interface \(q(x)\in B_0\).  We write
\[
x:p(x)\leadsto q(x).
\]

\begin{definition}[Product-state Mealy morphism]
\label{def:product-state-mealy-morphism}
A \emph{product-state Mealy morphism}
\[
(P,\delta,\omega):\mathbb A\rightsquigarrow\mathbb B
\]
is a Mealy morphism equipped with operations
\[
\top_P:1_{\mathcal E}\to P
\]
and
\[
\otimes_P:P\times P\to P
\]
such that
\[
p\top_P=\top_{\mathbb A},
\qquad
q\top_P=\top_{\mathbb B},
\]
and
\[
p(x\otimes_P y)=p(x)\times_{\mathbb A}p(y),
\qquad
q(x\otimes_P y)=q(x)\times_{\mathbb B}q(y).
\]
Equivalently, as morphisms \(P\times P\to A_0\) and \(P\times P\to B_0\),
\[
p\otimes_P
=
\times_{\mathbb A}\circ(p\times p),
\qquad
q\otimes_P
=
\times_{\mathbb B}\circ(q\times q).
\]
\end{definition}

Thus product-state structure is chosen structure on the state object lying
over the chosen product objects of the input and output internal categories.
It is object-level product structure on states.  It does not, by itself,
assert preservation of terminal arrows, projections, or pairings.

If
\[
x:A\leadsto B
\qquad\text{and}\qquad
y:X\leadsto Y,
\]
then
\[
x\otimes_P y:A\times X\leadsto B\times Y.
\]

This is the amount of product structure needed to type the Conway
loop-decomposition law below.  Under the target-to-source Mealy convention,
one should not expect the usual product-preservation equations for projections
and pairings to have the same shape as they do for functors.  The Conway laws
below isolate the product behavior relevant to recursive systems.

\section{Conway Mealy morphisms}
\label{sec:conway-mealy-morphisms}

Let
\[
\mathbb A
\qquad\text{and}\qquad
\mathbb B
\]
be internal Conway categories in \(\mathcal E\).  Let
\[
(P,\delta,\omega):\mathbb A\rightsquigarrow\mathbb B
\]
be a product-state Mealy morphism.

Because the Mealy convention is target-to-source, a recursive arrow
\[
f:A\times X\to X
\]
is processed at a state lying over \(X\).  If
\[
x:X\leadsto Y,
\]
then \(f\) is admissible at \(x\), and the transition gives
\[
\delta(f,x):A\times X\leadsto q(\delta(f,x))
\]
together with an output arrow
\[
\omega(f,x):q(\delta(f,x))\to Y.
\]

For this output to be a recursive equation in \(\mathbb B\), the updated state
must decompose as the product of a parameter state and the original recursive
state.  We now formulate this internally.

Let
\[
D_{\mathbb A}
\]
be the dagger-input context of \(\mathbb A\).  Thus \(D_{\mathbb A}\) has
projections, denoted in generalized-element notation by
\[
A,X:D_{\mathbb A}\to A_0,
\qquad
f:D_{\mathbb A}\to A_1,
\]
where
\[
f:A\times X\to X.
\]
Form the pullback
\[
R_P
:=
D_{\mathbb A}\times_{X,A_0,p}P.
\]
A generalized element of \(R_P\) is a tuple
\[
(A,X,f;x)
\]
with
\[
f:A\times X\to X,
\qquad
p(x)=X.
\]

\begin{definition}[Loop decomposition]
\label{def:loop-decomposition}
A product-state Mealy morphism
\[
(P,\delta,\omega):\mathbb A\rightsquigarrow\mathbb B
\]
between internal Conway categories has \emph{loop decomposition} if it is
equipped with a morphism
\[
\lambda:R_P\to P
\]
such that
\[
p\lambda=A
\]
and
\[
\delta(f,x)=\lambda(f,x)\otimes_P x
\]
as morphisms
\[
R_P\to P.
\]
In generalized-element notation, for every recursive input
\[
f:A\times X\to X
\]
and every state
\[
x:X\leadsto Y,
\]
we write
\[
\lambda(f,x):A\leadsto B
\]
for the parameter state determined by \(\lambda\), where
\[
B:=q(\lambda(f,x)).
\]
\end{definition}

When loop decomposition holds, the output arrow
\[
\omega(f,x)
\]
has type
\[
\omega(f,x):
q(\lambda(f,x))\times q(x)\to q(x).
\]
Indeed,
\[
\delta(f,x)=\lambda(f,x)\otimes_P x
\]
implies
\[
q(\delta(f,x))
=
q(\lambda(f,x))\times q(x),
\]
and the Mealy output typing gives
\[
\omega(f,x):q(\delta(f,x))\to q(x).
\]
Thus \(\omega(f,x)\) is itself a recursive arrow in the output Conway
category \(\mathbb B\).  We write
\[
f^{\omega,x}
:=
\omega(f,x)
:
q(\lambda(f,x))\times q(x)\to q(x).
\]

Formally, loop decomposition induces a map
\[
R_P\to D_{\mathbb B}
\]
sending
\[
(A,X,f;x)
\]
to
\[
\bigl(q(\lambda(f,x)),q(x),\omega(f,x)\bigr).
\]
The expression
\[
\bigl(f^{\omega,x}\bigr)^\dagger
\]
means the composite of this map \(R_P\to D_{\mathbb B}\) with the dagger
operation of \(\mathbb B\).

The solution
\[
f^\dagger:A\to X
\]
is also admissible at \(x:X\leadsto Y\), and processing it gives
\[
\delta(f^\dagger,x):A\leadsto q(\delta(f^\dagger,x)),
\]
together with
\[
\omega(f^\dagger,x):
q(\delta(f^\dagger,x))\to Y.
\]

\begin{definition}[Conway Mealy morphism]
\label{def:conway-mealy-morphism}
A \emph{Conway Mealy morphism}
\[
(P,\delta,\omega):\mathbb A\rightsquigarrow\mathbb B
\]
between internal Conway categories is a product-state Mealy morphism equipped
with loop decomposition
\[
\lambda:R_P\to P
\]
and satisfying the dagger preservation laws
\[
\delta(f^\dagger,x)=\lambda(f,x)
\]
as morphisms
\[
R_P\to P,
\]
and
\[
\omega(f^\dagger,x)
=
\bigl(f^{\omega,x}\bigr)^\dagger
\]
as morphisms
\[
R_P\to B_1.
\]
Equivalently, if processing the recursive equation
\[
f:A\times X\to X
\]
at the recursive state
\[
x:X\leadsto Y
\]
gives the output recursive equation
\[
f^{\omega,x}:B\times Y\to Y,
\]
then processing the chosen solution
\[
f^\dagger:A\to X
\]
at \(x\) gives the chosen solution
\[
(f^{\omega,x})^\dagger:B\to Y.
\]
\end{definition}

This is the target-to-source form of Conway preservation appropriate to the
Mealy convention
\[
\Phi:P\circ A\Rightarrow B\circ P.
\]

\section{Operational semantics of Conway Mealy morphisms}
\label{sec:operational-semantics-conway-mealy}

The ordinary Mealy calculus of the previous chapter describes a typed state
machine.  The Conway refinement describes what such a state machine does to
recursive systems and their chosen solutions.

Let
\[
(P,\delta,\omega):\mathbb A\rightsquigarrow\mathbb B
\]
be a Conway Mealy morphism.  A state
\[
x:X\leadsto Y
\]
relates an input object \(X\) of \(\mathbb A\) to an output object \(Y\) of
\(\mathbb B\).

A recursive equation in \(\mathbb A\) has the form
\[
f:A\times X\to X.
\]
Operationally, \(X\) is the recursive variable object and \(A\) is the
parameter object.  Since the Mealy machine processes arrows target-to-source,
the equation \(f\) is admissible at any state
\[
x:X\leadsto Y.
\]

Processing \(f\) at \(x\) gives
\[
x
\xrightarrow{\;f/\omega(f,x)\;}
\delta(f,x).
\]
The loop-decomposition law says that the updated state factors as
\[
\delta(f,x)=\lambda(f,x)\otimes_P x.
\]
Thus the machine separates the updated state into a parameter component and
the original recursive component:
\[
\lambda(f,x):A\leadsto B,
\qquad
x:X\leadsto Y.
\]
The emitted output arrow is therefore typed as
\[
\omega(f,x):B\times Y\to Y,
\]
where
\[
B=q(\lambda(f,x)),
\qquad
Y=q(x).
\]
This emitted arrow is a recursive equation in the output Conway category.

Hence a Conway Mealy morphism transforms the recursive equation
\[
f:A\times X\to X
\]
at the state \(x:X\leadsto Y\) into the recursive equation
\[
f^{\omega,x}:B\times Y\to Y.
\]

Now consider the chosen solution
\[
f^\dagger:A\to X.
\]
This arrow is again admissible at the recursive state \(x:X\leadsto Y\).  The
Conway Mealy law says that processing the solution is the same as solving the
processed equation:
\[
x
\xrightarrow{\;f^\dagger/(f^{\omega,x})^\dagger\;}
\lambda(f,x).
\]
Thus the state changes from the recursive state \(x\) to the parameter state
\(\lambda(f,x)\), and the emitted arrow is the chosen solution
\[
(f^{\omega,x})^\dagger:B\to Y.
\]

In slogan form:
\[
\boxed{
\text{run the recursive system, then solve}
=
\text{solve, then run the solution}.
}
\]
The equality is not an equality of untyped functions; it is the pair of typed
equations
\[
\delta(f^\dagger,x)=\lambda(f,x),
\qquad
\omega(f^\dagger,x)=\bigl(f^{\omega,x}\bigr)^\dagger.
\]

\subsection{Natural-deduction style rules}
\label{subsec:natural-deduction-conway-mealy}

The typing of the Conway transition can be displayed as rules.

The state formation rule is
\[
\frac{x\in P}
     {x:p(x)\leadsto q(x)}.
\]

The recursive-input rule is
\[
\frac{
f:A\times X\to_{\mathbb A}X
\qquad
x:X\leadsto Y
}
{
f\ \text{is admissible at}\ x
}.
\]

The Conway execution rule is
\[
\frac{
f:A\times X\to_{\mathbb A}X
\qquad
x:X\leadsto Y
}
{
x
\xrightarrow{\;f/f^{\omega,x}\;}
\lambda(f,x)\otimes_P x
}.
\]
The conclusion carries the typings
\[
\lambda(f,x):A\leadsto B,
\]
and
\[
f^{\omega,x}:B\times Y\to_{\mathbb B}Y,
\]
where
\[
B=q(\lambda(f,x)).
\]

The solution rule is
\[
\frac{
f:A\times X\to_{\mathbb A}X
\qquad
x:X\leadsto Y
}
{
x
\xrightarrow{\;f^\dagger/(f^{\omega,x})^\dagger\;}
\lambda(f,x)
}.
\]
This is exactly the dagger preservation law:
\[
\delta(f^\dagger,x)=\lambda(f,x),
\qquad
\omega(f^\dagger,x)=\bigl(f^{\omega,x}\bigr)^\dagger.
\]

\subsection{Interpretation}
\label{subsec:interpretation-conway-operational}

The role of loop decomposition is to make the output of a recursive input
again recursive.  Without it, processing
\[
f:A\times X\to X
\]
at a state \(x:X\leadsto Y\) would produce an output arrow
\[
\omega(f,x):q(\delta(f,x))\to Y.
\]
There would be no reason for the source object \(q(\delta(f,x))\) to have the
form
\[
B\times Y.
\]
Loop decomposition supplies precisely this factorization by requiring
\[
\delta(f,x)=\lambda(f,x)\otimes_P x.
\]

Thus a Conway Mealy morphism is a typed state machine that sends recursive
systems to recursive systems and chosen solutions to chosen solutions.

\section{Composition of Conway Mealy morphisms}
\label{sec:composition-conway-mealy}

Conway Mealy morphisms are closed under the pipeline composition of the
previous chapter.

Let
\[
(P,\delta_P,\omega_P):\mathbb A\rightsquigarrow\mathbb B
\]
and
\[
(Q,\delta_Q,\omega_Q):\mathbb B\rightsquigarrow\mathbb C
\]
be Conway Mealy morphisms.  Their composite has state object
\[
P\times_{q_P,B_0,p_Q}Q.
\]
A composite state is written
\[
(x,y)
\]
with
\[
q_P(x)=p_Q(y).
\]
The input and output interfaces are
\[
p_{QP}(x,y)=p_P(x),
\qquad
q_{QP}(x,y)=q_Q(y).
\]

The product-state structure is defined componentwise:
\[
\top_{QP}:=(\top_P,\top_Q),
\]
and
\[
(x,y)\otimes_{QP}(x',y')
=
(x\otimes_P x',\,y\otimes_Q y').
\]
This is well typed because
\[
q_P(x\otimes_P x')
=
q_P(x)\times q_P(x')
=
p_Q(y)\times p_Q(y')
=
p_Q(y\otimes_Q y').
\]

Let
\[
f:A\times X\to X
\]
be a recursive arrow in \(\mathbb A\), and let
\[
(x,y):X\leadsto Z
\]
be a composite recursive state.  First process \(f\) through \(P\).  Since
\(P\) is Conway Mealy, there is a parameter state
\[
a=\lambda_P(f,x):A\leadsto B
\]
such that
\[
\delta_P(f,x)=a\otimes_P x.
\]
The output recursive equation is
\[
g:=f^{\omega_P,x}
=
\omega_P(f,x):
B\times q_P(x)\to q_P(x).
\]

Now process \(g\) through \(Q\) at the state \(y\).  Since \(Q\) is Conway
Mealy, there is a parameter state
\[
b=\lambda_Q(g,y):B\leadsto C
\]
such that
\[
\delta_Q(g,y)=b\otimes_Q y.
\]
The output recursive equation is
\[
h:=g^{\omega_Q,y}
=
\omega_Q(g,y):
C\times q_Q(y)\to q_Q(y).
\]

Because
\[
q_P(a)=B=p_Q(b),
\]
the pair \((a,b)\) is a well-typed state of the composite.  Define the loop
decomposition for the composite by
\[
\lambda_{QP}(f,(x,y))
:=
(a,b).
\]
Then
\[
\delta_{QP}(f,(x,y))
=
\bigl(\delta_P(f,x),\,\delta_Q(g,y)\bigr)
=
(a\otimes_P x,\,b\otimes_Q y)
=
(a,b)\otimes_{QP}(x,y).
\]

For dagger preservation, the Conway law for \(P\) gives
\[
\delta_P(f^\dagger,x)=a,
\qquad
\omega_P(f^\dagger,x)=g^\dagger.
\]
Then the Conway law for \(Q\) gives
\[
\delta_Q(g^\dagger,y)=b,
\qquad
\omega_Q(g^\dagger,y)=h^\dagger.
\]
Therefore the composite satisfies
\[
\delta_{QP}(f^\dagger,(x,y))
=
(a,b)
=
\lambda_{QP}(f,(x,y)),
\]
and
\[
\omega_{QP}(f^\dagger,(x,y))
=
h^\dagger.
\]

\begin{proposition}
\label{prop:conway-mealy-closed-composition}
Conway Mealy morphisms are closed under Mealy composition.  The identity
Mealy morphism on an internal Conway category is a Conway Mealy morphism.
\end{proposition}

\begin{proof}
The preceding calculation proves closure under composition.  For the identity
Mealy morphism on \(\mathbb A\), the state object is \(A_0\), the state product
is the chosen product of objects in \(\mathbb A\), and processing an arrow
emits the same arrow.

For a recursive arrow
\[
f:A\times X\to X
\]
at the state \(X\), the transition updates to the state
\[
A\times X,
\]
which is the product of the parameter state \(A\) and the recursive state
\(X\).  The emitted recursive equation is \(f\) itself.  Thus
\[
\lambda(f,X)=A.
\]
Processing
\[
f^\dagger:A\to X
\]
at \(X\) updates to \(A\) and emits \(f^\dagger\).  Hence the identity is
Conway Mealy.
\end{proof}

\section{Conway Mealy transformations}
\label{sec:conway-mealy-transformations}

Let
\[
(P,\delta_P,\omega_P)
\qquad\text{and}\qquad
(Q,\delta_Q,\omega_Q)
\]
be Conway Mealy morphisms
\[
\mathbb A\rightsquigarrow\mathbb B.
\]

\begin{definition}[Conway Mealy transformation]
\label{def:conway-mealy-transformation}
A \emph{Conway Mealy transformation}
\[
\theta:P\Rightarrow Q
\]
is a Mealy transformation
\[
\theta:P\to Q
\]
which preserves the product-state structure:
\[
\theta\top_P=\top_Q,
\qquad
\theta(x\otimes_P y)=\theta(x)\otimes_Q\theta(y),
\]
and preserves loop decomposition:
\[
\theta(\lambda_P(f,x))
=
\lambda_Q(f,\theta x).
\]
Formally, the last equation is an equality of morphisms out of the recursive
state context \(R_P\).
\end{definition}

The usual Mealy transformation equations already say
\[
\theta(\delta_P(a,x))
=
\delta_Q(a,\theta x),
\]
and
\[
\omega_P(a,x)
=
\omega_Q(a,\theta x).
\]
Together with preservation of \(\lambda\), these equations are compatible
with dagger preservation.

\begin{proposition}
\label{prop:conway-mealy-transformations-category}
For fixed internal Conway categories \(\mathbb A\) and \(\mathbb B\), Conway
Mealy morphisms
\[
\mathbb A\rightsquigarrow\mathbb B
\]
and Conway Mealy transformations between them form a category.
\end{proposition}

\begin{proof}
Identities and composites are inherited from Mealy transformations.  The
additional equations for \(\top\), \(\otimes\), and \(\lambda\) are preserved
by ordinary composition of maps in \(\mathcal E\).
\end{proof}

\begin{proposition}[Horizontal composition of Conway Mealy transformations]
\label{prop:horizontal-composition-conway-mealy-transformations}
Let
\[
\theta:P\Rightarrow P'
\]
be a Conway Mealy transformation
\[
\mathbb A\rightsquigarrow\mathbb B
\]
and let
\[
\kappa:Q\Rightarrow Q'
\]
be a Conway Mealy transformation
\[
\mathbb B\rightsquigarrow\mathbb C.
\]
Then the horizontal composite
\[
\theta\times\kappa:
P\times_{B_0}Q
\to
P'\times_{B_0}Q'
\]
defined by
\[
(x,y)\longmapsto(\theta x,\kappa y)
\]
is a Conway Mealy transformation.
\end{proposition}

\begin{proof}
The map is well defined because \(\theta\) and \(\kappa\) are maps of spans:
\[
q_{P'}\theta=q_P,
\qquad
p_{Q'}\kappa=p_Q.
\]
The ordinary Mealy transformation axioms are inherited from the horizontal
composition of 2-cells in \(\mathbf{Span}(\mathcal E)\).

It remains to check the extra Conway data.  Preservation of the terminal
state and product-state operations is componentwise:
\[
(\theta\times\kappa)(\top_P,\top_Q)
=
(\top_{P'},\top_{Q'}),
\]
and
\[
(\theta\times\kappa)
\bigl((x,y)\otimes(x',y')\bigr)
=
(\theta x,\kappa y)\otimes(\theta x',\kappa y').
\]

For loop decomposition, let
\[
f:A\times X\to X
\]
be a recursive arrow of \(\mathbb A\), and let
\[
(x,y):X\leadsto Z
\]
be a composite state.  Write
\[
g=\omega_P(f,x).
\]
Since \(\theta\) is a Mealy transformation,
\[
\omega_P(f,x)=\omega_{P'}(f,\theta x).
\]
Therefore the recursive equation fed into \(Q\) and \(Q'\) is the same.  Hence
\[
\begin{aligned}
(\theta\times\kappa)\lambda_{QP}(f,(x,y))
&=
(\theta\times\kappa)
\bigl(\lambda_P(f,x),\lambda_Q(g,y)\bigr)
\\
&=
\bigl(\lambda_{P'}(f,\theta x),\lambda_{Q'}(g,\kappa y)\bigr)
\\
&=
\lambda_{Q'P'}(f,(\theta x,\kappa y)).
\end{aligned}
\]
Thus the horizontal composite preserves loop decomposition.
\end{proof}

\section{Internal Conway functors as representable cases}
\label{sec:internal-conway-functors-representable}

The previous chapter showed that internal functors are exactly Mealy morphisms
carried by canonical representable spans.  We now record the Conway refinement
of that fact.

Let
\[
F:\mathbb A\to\mathbb B
\]
be an internal Conway functor.  Its object map
\[
F_0:A_0\to B_0
\]
determines the canonical representable span
\[
F_{0!}
=
\left(
A_0
\xleftarrow{1_{A_0}}
A_0
\xrightarrow{F_0}
B_0
\right).
\]
As in the previous chapter, the arrow map
\[
F_1:A_1\to B_1
\]
defines a Mealy structure on \(F_{0!}\).  The state object is \(A_0\), so a
state is just an object of \(\mathbb A\).  An arrow
\[
a:u\to v
\]
is processed at the state \(v\), updates the state to \(u\), and emits
\[
F_1(a):F_0u\to F_0v.
\]

The product-state structure is the chosen finite-product structure on
\(\mathbb A\):
\[
\top_F:=\top_{\mathbb A},
\qquad
X\otimes_F Y:=X\times_{\mathbb A}Y.
\]
Because \(F\) preserves chosen finite products strictly,
\[
F_0(X\times_{\mathbb A}Y)
=
F_0X\times_{\mathbb B}F_0Y,
\]
and
\[
F_0(\top_{\mathbb A})=\top_{\mathbb B}.
\]

Let
\[
f:A\times X\to X
\]
be a recursive arrow of \(\mathbb A\), and process it at the state \(X\).  The
updated state is
\[
A\times X
=
A\otimes_F X,
\]
so the loop-decomposition state is
\[
\lambda(f,X)=A.
\]
The emitted recursive equation is
\[
F_1(f):F_0(A\times X)\to F_0X.
\]
Since \(F\) preserves products, this is an arrow
\[
F_0A\times F_0X\to F_0X.
\]

Processing the solution
\[
f^\dagger:A\to X
\]
at the state \(X\) emits
\[
F_1(f^\dagger):F_0A\to F_0X.
\]
Since \(F\) preserves the Conway dagger,
\[
F_1(f^\dagger)
=
F_1(f)^\dagger.
\]
Thus the representable Mealy morphism induced by \(F\) is Conway Mealy.

\begin{proposition}
\label{prop:internal-conway-functors-representable-conway-mealy}
Every internal Conway functor determines a representable Conway Mealy
morphism.
\end{proposition}

\begin{proof}
The previous chapter identifies the internal functor \(F\) with the Mealy
morphism carried by the canonical representable span \(F_{0!}\).  Preservation
of finite products gives the product-state structure described above.
Preservation of the dagger gives the Conway Mealy dagger law.
\end{proof}

Conversely, a Conway Mealy morphism carried by a canonical representable span
determines an internal functor by the theorem of the previous chapter.  The
product-state equations force this induced functor to preserve the chosen
terminal object and binary product objects strictly, but they do not by
themselves force preservation of terminal arrows, projections, or pairings.
Hence preservation of the full chosen finite-product structure must be assumed
separately.

\begin{proposition}
\label{prop:representable-conway-mealy-to-internal-conway-functor}
Let
\[
(P,\delta,\omega):\mathbb A\rightsquigarrow\mathbb B
\]
be a Conway Mealy morphism carried by a canonical representable span
\[
F_{0!}
=
\left(
A_0
\xleftarrow{1_{A_0}}
A_0
\xrightarrow{F_0}
B_0
\right).
\]
Let
\[
F:\mathbb A\to\mathbb B
\]
be the internal functor induced by the underlying representable Mealy
morphism.  If \(F\) preserves the chosen finite-product structure, then \(F\)
is an internal Conway functor.
\end{proposition}

\begin{proof}
It remains only to check dagger preservation.  Let
\[
f:A\times X\to X
\]
be a recursive arrow in \(\mathbb A\).  In the representable carrier, the
state over \(X\) is \(X\) itself.  Loop decomposition identifies the parameter
state as \(A\), and the emitted recursive equation is \(F_1(f)\), regarded
using the assumed product preservation as a morphism
\[
F_0A\times F_0X\to F_0X.
\]
The Conway Mealy dagger law gives
\[
F_1(f^\dagger)=F_1(f)^\dagger.
\]
Thus \(F\) preserves the Conway dagger.  Since it also preserves the chosen
finite products by assumption, it is an internal Conway functor.
\end{proof}

\section{The bicategory of internal Conway categories and Conway Mealy morphisms}
\label{sec:bicategory-internal-conway-mealy}

\begin{definition}[The Conway Mealy bicategory]
\label{def:conway-mealy-bicategory}
Let \(\mathcal E\) be finitely complete.  The bicategory
\[
\mathbf{CnvMealy}(\mathcal E)
\]
has:
\begin{itemize}
\item objects: internal Conway categories in \(\mathcal E\);
\item 1-cells: Conway Mealy morphisms;
\item 2-cells: Conway Mealy transformations.
\end{itemize}
Composition is the pipeline composition of Mealy morphisms, equipped with the
componentwise product-state structure and loop decomposition described in
Section~\ref{sec:composition-conway-mealy}.
\end{definition}

\begin{proposition}
\label{prop:cnvmealy-is-bicategory}
The data of Definition~\ref{def:conway-mealy-bicategory} form a bicategory.
\end{proposition}

\begin{proof}
For fixed objects \(\mathbb A,\mathbb B\), the hom-category is given by
Proposition~\ref{prop:conway-mealy-transformations-category}.  Horizontal
composition of 1-cells is given by pipeline composition, and closure under
this composition is Proposition~\ref{prop:conway-mealy-closed-composition}.
Horizontal composition of 2-cells is
Proposition~\ref{prop:horizontal-composition-conway-mealy-transformations}.

The associativity and unit constraints are inherited from the bicategory of
spans.  Explicitly, the associator identifies
\[
(P\times_{B_0}Q)\times_{C_0}R
\]
with
\[
P\times_{B_0}(Q\times_{C_0}R)
\]
by the canonical pullback comparison map, which in generalized-element
notation sends
\[
((x,y),z)
\]
to
\[
(x,(y,z)).
\]
Because product-state structures and loop decompositions for composites are
defined componentwise, this comparison preserves \(\top\), \(\otimes\), and
\(\lambda\).  Hence the associator is a Conway Mealy transformation.  The same
argument applies to the left and right unitors, which remove or insert the
identity state component.

The pentagon and triangle axioms follow from the corresponding axioms in
\(\mathbf{Span}(\mathcal E)\), since all additional Conway Mealy structure is
preserved by the canonical pullback comparison isomorphisms.
\end{proof}

\section{Conway span-monads with Mealy morphisms}
\label{sec:conway-span-monads-with-mealy-morphisms}

We can express the same bicategory through span-monads.

\begin{definition}[Conway span-monad Mealy morphism]
\label{def:conway-span-monad-mealy-morphism}
Let
\[
C
\qquad\text{and}\qquad
D
\]
be Conway span-monads corresponding to internal Conway categories
\[
\mathbb C
\qquad\text{and}\qquad
\mathbb D.
\]
A \emph{Conway span-monad Mealy morphism}
\[
C\rightsquigarrow D
\]
is a span-monad morphism
\[
(P,\Phi):C\rightsquigarrow D
\]
with structure cell
\[
\Phi:P\circ C\Rightarrow D\circ P
\]
whose associated Mealy morphism
\[
P:\mathbb C\rightsquigarrow\mathbb D
\]
is a Conway Mealy morphism.
\end{definition}

\begin{definition}[The bicategory \(\mathbf{CnvSpanMnd}_{\mathrm{Mealy}}(\mathcal E)\)]
\label{def:cnvspanmnd-mealy}
Let
\[
\mathbf{CnvSpanMnd}_{\mathrm{Mealy}}(\mathcal E)
\]
be the bicategory whose objects are Conway span-monads in
\[
\mathbf{Span}(\mathcal E),
\]
whose 1-cells are Conway span-monad Mealy morphisms, and whose 2-cells are
the corresponding Conway Mealy transformations.
\end{definition}

\begin{theorem}[Span-monadic form of the Conway Mealy calculus]
\label{thm:span-monadic-conway-mealy-calculus}
Let \(\mathcal E\) be finitely complete.  The correspondence between
span-monads and internal categories identifies
\[
\mathbf{CnvSpanMnd}_{\mathrm{Mealy}}(\mathcal E)
\]
with
\[
\mathbf{CnvMealy}(\mathcal E).
\]
\end{theorem}

\begin{proof}
On objects, this is
Theorem~\ref{thm:conway-span-monads-internal-conway-categories}.  On 1-cells,
it is the Mealy correspondence from the previous chapter: a span-monad
morphism with structure cell
\[
\Phi:P\circ C\Rightarrow D\circ P
\]
is exactly a Mealy morphism between the associated internal categories.
The additional product-state, loop-decomposition, and dagger-preservation
data are the same on both sides.  The statement for 2-cells follows from the
definition of Conway Mealy transformation.  The compatibility with
composition, associators, and unitors is inherited from
\(\mathbf{Span}(\mathcal E)\) together with
Proposition~\ref{prop:cnvmealy-is-bicategory}.
\end{proof}

\section{The case \(\mathcal E=\mathbf{Set}\)}
\label{sec:conway-set-case}

When
\[
\mathcal E=\mathbf{Set},
\]
internal categories are small categories.  Hence internal Conway categories in
\(\mathbf{Set}\) are small Conway categories.

A Conway span-monad in
\[
\mathbf{Span}(\mathbf{Set})
\]
is therefore a small category with chosen finite products and a Conway dagger,
presented as a span-monad.

A Conway Mealy morphism between small Conway categories is a typed state
machine
\[
(P,\delta,\omega):\mathbb A\rightsquigarrow\mathbb B
\]
whose states have input and output interfaces
\[
p(x)\in A_0,
\qquad
q(x)\in B_0,
\]
whose transitions process input arrows target-to-source, and whose
product-state and loop-decomposition structure transports recursive systems
to recursive systems and chosen solutions to chosen solutions.

The ordinary Conway functors are exactly those representable Conway Mealy
morphisms whose induced functors preserve the chosen finite-product structure
and the Conway dagger.

\section{Main consolidation}
\label{sec:conway-main-consolidation}

The chapter can be summarized as follows.

\begin{theorem}[Conway span-monads and Conway Mealy morphisms]
\label{thm:conway-main-consolidation}
Let \(\mathcal E\) be finitely complete.

\begin{enumerate}
\item The finite-limit theory
\[
\ConvCatTh
\]
has as its models in \(\mathcal E\) precisely internal Conway categories in
\(\mathcal E\).

\item Conway span-monads in
\[
\mathbf{Span}(\mathcal E)
\]
are precisely span-monads whose associated internal categories are models of
\[
\ConvCatTh.
\]
Thus Conway span-monads are internal Conway categories presented through the
span-monad correspondence.

\item A span-monad morphism with structure cell
\[
\Phi:P\circ A\Rightarrow B\circ P
\]
is, under the notation of the previous chapter, a Mealy morphism
\[
(P,\delta,\omega):\mathbb A\rightsquigarrow\mathbb B.
\]

\item A Conway Mealy morphism is a Mealy morphism between internal Conway
categories equipped with product-state structure, loop decomposition, and
dagger preservation.

\item Internal Conway functors determine representable Conway Mealy
morphisms.  Conversely, a representable Conway Mealy morphism whose induced
internal functor preserves the chosen finite-product structure is an internal
Conway functor.

\item The span-monadic presentation and the internal-category presentation
give the same Conway Mealy calculus:
\[
\mathbf{CnvSpanMnd}_{\mathrm{Mealy}}(\mathcal E)
\simeq
\mathbf{CnvMealy}(\mathcal E).
\]
\end{enumerate}
\end{theorem}

\begin{proof}
The first assertion is Proposition~\ref{prop:models-convcatth}.  The second
is Theorem~\ref{thm:conway-span-monads-internal-conway-categories}.  The
third is the Mealy-morphism theorem of the previous chapter.  The fourth is
Definition~\ref{def:conway-mealy-morphism}.  The fifth is
Propositions~\ref{prop:internal-conway-functors-representable-conway-mealy}
and
\ref{prop:representable-conway-mealy-to-internal-conway-functor}.  The sixth
is Theorem~\ref{thm:span-monadic-conway-mealy-calculus}.
\end{proof}

The conceptual picture is therefore:
\[
\boxed{
\text{Lawvere theories}
\longleftrightarrow
\text{finitary monads}
}
\]
in the ordinary algebraic setting, and
\[
\boxed{
\text{finite-limit theories of internal categories}
\longleftrightarrow
\text{span-monads}
}
\]
in the internal setting.  Adding Conway structure on the theory side gives
Conway categories; transporting it through the span-monad correspondence gives
Conway span-monads.  The morphisms native to spans are Mealy morphisms, and
their Conway refinement is the stateful calculus described above.